% Appendix A of Deep Learning as an Experimental Science: the statistical toolkit.
\labappendix{A Statistical Toolkit}{The statistics the experiments share}{ch:toolkit}
\begin{refsection}[chapters/17-toolkit.bib]

\begin{chaptersummary}
A handful of statistical tools recur across the chapters of this book: intervals for a mean and for a difference of means; the bootstrap and the permutation test; paired comparisons of two classifiers on the same test items; corrections for many comparisons; the distribution of the best of several draws; the delta method; nonlinear least squares; variance components and factorial designs; sample sizes; rank statistics; and intervals for proportions near zero. This appendix states what the chapters use and proves what can be proved in a page; for the rest it gives a sketch and a reference.
\end{chaptersummary}

The appendix is a reference, not a chapter to be read in order. Intervals and tests for means (\cref{sec:toolkit:mean}) are used everywhere; resampling and paired tests (\cref{sec:toolkit:resampling}) in \cref{ch:seeds,ch:instability,ch:aggregates,ch:budgets,ch:fitting,ch:protocol,ch:extrapolation,ch:instruments}; multiple comparisons (\cref{sec:toolkit:multiple}) in \cref{ch:aggregates}; order statistics (\cref{sec:toolkit:orderstat}) in \cref{ch:seeds,ch:budgets,ch:controls,ch:emergence}; the delta method and nonlinear least squares (\cref{sec:toolkit:delta,sec:toolkit:nls}) in \cref{ch:fitting,ch:protocol,ch:extrapolation,ch:emergence}; variance components and factorial designs (\cref{sec:toolkit:design}) in \cref{ch:seeds,ch:controls,ch:confounding,ch:ablations,ch:sharpness,ch:instruments}; sample sizes (\cref{sec:toolkit:power}) in the epilogue and in \cref{ch:seeds,ch:instability,ch:controls}; rank statistics (\cref{sec:toolkit:ranks}) in \cref{ch:aggregates,ch:controls,ch:sharpness,ch:instruments}; and binomial intervals (\cref{sec:toolkit:binomial}) in \cref{ch:seeds,ch:extrapolation,ch:emergence}.

Throughout, $Y_1,\dots,Y_n$ are outcomes of runs, $\bar Y$ is their mean and $S^2=\frac1{n-1}\sum_i(Y_i-\bar Y)^2$ their sample variance, for $n\ge2$. As in the rest of the book, $\Phi$ is the standard normal distribution function, $z_q$ its $q$-quantile, and $t_{\nu,q}$ the $q$-quantile of Student's $t$ distribution with $\nu$ degrees of freedom. A $p$-value for a true null hypothesis is called \term{valid}\index{p-value!valid} when $\Prob(p\le u)\le u$ for every $u\in[0,1]$.

\section{Intervals for means}\label{sec:toolkit:mean}

If $Y_1,\dots,Y_n$ are independent with mean $\mu$ and variance $\sigma^2$, then $\E\bar Y=\mu$, $\Var\bar Y=\sigma^2/n$, and $\E S^2=\sigma^2$, so that $S/\sqrt n$, the \term{standard error}\index{standard error}, estimates the standard deviation of $\bar Y$.

\begin{theorem}[$t$-intervals and Welch's test]\label{thm:toolkit:tinterval}\index{t-interval}\index{Welch's test}
\begin{enumerate}
\item If $Y_1,\dots,Y_n$ are independent and normal with mean $\mu$ and variance $\sigma^2>0$, then $(\bar Y-\mu)/(S/\sqrt n)$ has Student's $t$ distribution with $n-1$ degrees of freedom, and the interval $\bar Y\pm t_{n-1,1-\alpha/2}\,S/\sqrt n$ contains $\mu$ with probability exactly $1-\alpha$.
\item If $Y_1,\dots,Y_n$ are independent and identically distributed with mean $\mu$ and finite positive variance, then the coverage of the same interval tends to $1-\alpha$ as $n\to\infty$.
\item For two independent samples, each independent and identically distributed, with sizes $n_1,n_2$, means $\mu_1,\mu_2$, and finite positive variances $\sigma_1^2,\sigma_2^2$, Welch's statistic
\[
W=\frac{\bar Y_1-\bar Y_2-(\mu_1-\mu_2)}{\sqrt{S_1^2/n_1+S_2^2/n_2}}
\]
converges in distribution to the standard normal as $n_1,n_2\to\infty$. For normal samples $W$ is approximately $t$-distributed with
$\nu=\bigl(S_1^2/n_1+S_2^2/n_2\bigr)^2\big/\bigl[(S_1^2/n_1)^2/(n_1-1)+(S_2^2/n_2)^2/(n_2-1)\bigr]$
degrees of freedom, and Welch's test of $\mu_1=\mu_2$ compares $W$, computed with $\mu_1-\mu_2=0$, with the quantiles of this distribution.
\end{enumerate}
\end{theorem}

\begin{proofsketch}
Part (i) rests on two facts about normal samples: $\bar Y$ and $S^2$ are independent, and $(n-1)S^2/\sigma^2$ has the $\chi^2$ distribution with $n-1$ degrees of freedom \parencite[Section~5.3]{CasellaBerger2002}. A standard normal divided by the square root of an independent $\chi^2_\nu/\nu$ has, by definition, Student's distribution with $\nu$ degrees of freedom \parencite{Student1908}. Part (ii) follows from the central limit theorem for $\sqrt n(\bar Y-\mu)/\sigma$, the convergence of $S$ to $\sigma$ in probability, Slutsky's lemma, and $t_{n-1,1-\alpha/2}\to z_{1-\alpha/2}$. Part (iii) is proved the same way for the asymptotic statement; the degrees of freedom come from matching the first two moments of $S_1^2/n_1+S_2^2/n_2$ to those of a scaled $\chi^2$ variable \parencite{Satterthwaite1946,Welch1947}, which makes the $t$ distribution an approximation rather than an exact law.
\end{proofsketch}

When the runs of two arms come in pairs that share a seed, a data order, or a set of test items, the comparison should use the differences $D_i=Y_{1i}-Y_{2i}$ and the one-sample interval of part (i) for $\E D$. Since $\Var D=\sigma_1^2+\sigma_2^2-2\Cov(Y_1,Y_2)$, pairing lowers the standard error whenever the paired outcomes are positively correlated, which is the usual case. It costs degrees of freedom, $n-1$ instead of up to $2n-2$, which matters only for very few runs or a weak correlation.

\section{Resampling}\label{sec:toolkit:resampling}

\begin{definition}[Bootstrap intervals]\label{def:toolkit:bootstrap}\index{bootstrap}
Let $\hat\theta=\hat\theta(Y_1,\dots,Y_n)$ be a statistic. A \term{bootstrap sample} draws $n$ values uniformly with replacement from $Y_1,\dots,Y_n$, and $\hat\theta$ computed on it is a \term{bootstrap replicate} $\hat\theta^*$. The \term{percentile interval}\index{bootstrap!percentile interval} at level $1-\alpha$ runs from the $\alpha/2$-quantile to the $(1-\alpha/2)$-quantile of the replicates computed on many independent bootstrap samples. When the data fall into strata, for instance tasks with several runs each, the \term{stratified bootstrap}\index{bootstrap!stratified} resamples with replacement within each stratum and keeps the strata fixed. The \term{parametric bootstrap} instead draws each new data set from a model fitted to the data, for instance by adding fresh normal errors of the estimated variance to a fitted curve, and recomputes $\hat\theta$ on it.
\end{definition}

\begin{theorem}[Consistency of the bootstrap]\label{thm:toolkit:bootstrap}
Let $Y_1,\dots,Y_n$ be independent and identically distributed with mean $\mu$ and finite positive variance $\sigma^2$, and let $\hat\theta=g(\bar Y)$ for a function $g$ that is continuously differentiable in a neighborhood of $\mu$ with $g'(\mu)\ne0$. Then, for almost every sequence of data, the conditional distribution of $\sqrt n(\hat\theta^*-\hat\theta)$ given $Y_1,\dots,Y_n$ converges to $\cN(0,g'(\mu)^2\sigma^2)$, which is also the limit of the distribution of $\sqrt n(\hat\theta-g(\mu))$. Consequently the percentile interval, with the quantiles taken of this conditional distribution (the limit of infinitely many replicates), has coverage tending to $1-\alpha$.
\end{theorem}

\begin{proofsketch}
\Textcite{BickelFreedman1981} prove the conditional central limit theorem for the bootstrap mean: for almost every sequence of data, $\sqrt n(\bar Y^*-\bar Y)$ converges conditionally to $\cN(0,\sigma^2)$. By the mean value theorem $\sqrt n\bigl(g(\bar Y^*)-g(\bar Y)\bigr)=g'(\xi)\sqrt n(\bar Y^*-\bar Y)$ for a point $\xi$ between the two means, and $g'(\xi)\to g'(\mu)$ because both means converge to $\mu$; this is the delta method (\cref{thm:toolkit:delta}) for the bootstrap. The conditional $q$-quantile $c^*_q$ of $\sqrt n(\hat\theta^*-\hat\theta)$ then tends to $\abs{g'(\mu)}\sigma z_q$. The interval contains $g(\mu)$ exactly when $-c^*_{1-\alpha/2}\le\sqrt n\bigl(\hat\theta-g(\mu)\bigr)\le-c^*_{\alpha/2}$, and the symmetry $z_{\alpha/2}=-z_{1-\alpha/2}$ of the normal limit makes the probability of this event tend to $1-\alpha$ \parencite[Chapter~23]{vanderVaart1998}.
\end{proofsketch}

With few runs the percentile interval tends to be too narrow: the bootstrap variance of a mean of $n$ runs is $\frac{n-1}n\,S^2/n$, and its percentiles, like normal quantiles, ignore the error in the estimated spread that $t$ quantiles allow for. With few runs in every stratum the stratified bootstrap keeps this shrinkage however many strata there are (\cref{ch:aggregates}); \textcite{BickelFreedman1984} study the bootstrap in stratified sampling. The bias-corrected and accelerated interval of \textcite{Efron1987} corrects for bias and skewness, and \textcite{EfronTibshirani1993} is the standard introduction. Under further regularity conditions consistency extends to smooth functionals of the empirical distribution and to the parametric bootstrap \parencite[Chapter~23]{vanderVaart1998}. The chapters also bootstrap fitted exponents, a case the theorem does not cover; there a simulated control, as in \cref{ch:extrapolation}, measures the coverage actually achieved.

A permutation test needs no asymptotics. Its null hypothesis is a symmetry: for two arms of runs, that the arm labels are exchangeable, as they are when the two arms have the same run distribution and the runs are independent.

\begin{theorem}[Permutation tests are exact]\label{thm:toolkit:permutation}\index{permutation test}
Let $G$ be a finite group acting on the data, and suppose that under the null hypothesis $gY$ has the same distribution as $Y$ for every $g\in G$. For a statistic $T$ put $p(Y)=\abs{G}^{-1}\,\#\{h\in G:T(hY)\ge T(Y)\}$. Then, under the null hypothesis, $\Prob(p(Y)\le u)\le u$ for every $u\in[0,1]$.
\end{theorem}

\begin{proof}
The plan is to show that in every orbit of the group at most a fraction $u$ of the points have a $p$-value at most $u$, and then to use the symmetry to turn this count into a probability.

Fix a data set $y$ and let $A=\{g\in G:p(gy)\le u\}$. Since $h\mapsto hg$ is a bijection of $G$, the $p$-value of $gy$ counts the same values $T(hy)$ as that of $y$, against a different threshold:
\[
p(gy)=\abs{G}^{-1}\,\#\{h\in G:T(hgy)\ge T(gy)\}=\abs{G}^{-1}\,\#\{h\in G:T(hy)\ge T(gy)\}.
\]
If $A$ is empty, then $\abs A=0\le u\abs G$. Otherwise choose $g^\star\in A$ for which $T(g^\star y)$ is smallest among the elements of $A$. Every $g\in A$ satisfies $T(gy)\ge T(g^\star y)$, so $A$ is contained in $\{h:T(hy)\ge T(g^\star y)\}$, a set with $\abs G\,p(g^\star y)\le u\abs G$ elements by the display. Hence $\abs{A}\le u\abs{G}$ for every $y$.

Now let the data be random. By the null hypothesis $p(gY)$ has the distribution of $p(Y)$ for each $g$, so averaging over $g$ and then applying the bound on $\abs A$ with $y=Y$ gives
\[
\Prob(p(Y)\le u)=\frac1{\abs G}\sum_{g\in G}\Prob(p(gY)\le u)=\E\Bigl[\frac{\#\{g:p(gY)\le u\}}{\abs G}\Bigr]\le u.\qedhere
\]
\end{proof}

For two arms of $n$ runs, $G$ is the group of permutations of the $2n$ pooled outcomes, acting by reassigning arm labels, and $T$ is for instance the difference of the arm means, or its absolute value for a two-sided test. When $\abs G$ is too large to enumerate, $M$ permutations drawn independently and uniformly from $G$, with statistics $T_1,\dots,T_M$, give the valid $p$-value $(1+\#\{b:T_b\ge T_{\mathrm{obs}}\})/(M+1)$: under the null hypothesis the observed statistic and the $M$ permuted ones are exchangeable, so the rank of the observed one among them is uniform up to ties. Dropping the two ones, as is often done, can produce a $p$-value of zero and is not valid \parencite{PhipsonSmyth2010}.

Two classifiers evaluated on the same test items are a paired design. Its analysis goes back to \textcite{McNemar1947}; \textcite{Dietterich1998} recommended McNemar's test for comparing two learning algorithms when each can be trained only once.

\begin{proposition}[Two classifiers on the same items]\label{thm:toolkit:mcnemar}\index{McNemar's test}
Two fixed classifiers are evaluated on $m$ test items drawn independently. Let $b$ count the items that the first classifies correctly and the second incorrectly, and $c$ the items with the reverse, and let $\pi_b,\pi_c$ be the probabilities of these two discordances.
\begin{enumerate}
\item If $\pi_b=\pi_c$, then conditionally on $b+c=s$ the count $b$ has the binomial distribution $\Bin(s,1/2)$.
\item The difference of the two accuracies is $(b-c)/m$, with variance $\bigl(\pi_b+\pi_c-(\pi_b-\pi_c)^2\bigr)/m$.
\end{enumerate}
\end{proposition}

\begin{proof}
Each item falls, independently of the others, into one of four categories: both correct, only the first correct, only the second correct, both wrong. Conditionally on which items are discordant, each discordant item is of the first kind with probability $\pi_b/(\pi_b+\pi_c)$, independently of the others; under $\pi_b=\pi_c$ this probability is $1/2$, which proves (i). For (ii), write the difference of accuracies as $\frac1m\sum_iX_i$ with $X_i=1$ for a discordance of the first kind, $-1$ for one of the second kind, and $0$ otherwise. Then $\E X_i=\pi_b-\pi_c$ and $\E X_i^2=\pi_b+\pi_c$, so $\Var X_i=\pi_b+\pi_c-(\pi_b-\pi_c)^2$, and the mean of $m$ independent copies has $1/m$ of this variance.
\end{proof}

McNemar's test of $\pi_b=\pi_c$ refers $b$ to the binomial distribution of part (i), or, when $b+c$ is large, $(b-c)^2/(b+c)$ to the $\chi^2$ distribution with one degree of freedom. Two classifiers of similar accuracy usually agree on most items, so $\pi_b+\pi_c$ is small and the paired variance in (ii) is far below $\bigl(p_1(1-p_1)+p_2(1-p_2)\bigr)/m$, the variance for accuracies $p_1,p_2$ measured on two independent test sets of $m$ items each.

\section{Many comparisons}\label{sec:toolkit:multiple}

Let $H_1,\dots,H_m$ be null hypotheses with valid $p$-values $p_1,\dots,p_m$; in this section $m$ counts hypotheses, not test items. Let $I_0$ be the set of true nulls and $m_0=\abs{I_0}$. A procedure rejects some of the hypotheses; let $R$ be the number of rejections and $V$ the number of true nulls among them. The \term{family-wise error rate}\index{family-wise error rate} is $\Prob(V\ge1)$, and the \term{false discovery rate}\index{false discovery rate} is $\E[V/\max(R,1)]$.

\begin{proposition}[Bonferroni]\label{thm:toolkit:bonferroni}
Rejecting $H_i$ when $p_i\le\alpha/m$ gives a family-wise error rate at most $m_0\alpha/m\le\alpha$, whatever the dependence between the $p$-values.
\end{proposition}

\begin{proof}
By the union bound and validity, $\Prob(V\ge1)\le\sum_{i\in I_0}\Prob(p_i\le\alpha/m)\le m_0\alpha/m$.
\end{proof}

\begin{theorem}[Holm's procedure]\label{thm:toolkit:holm}\index{Holm's procedure}
Order the $p$-values as $p_{(1)}\le\dots\le p_{(m)}$. Let $j$ be the smallest index with $p_{(j)}>\alpha/(m-j+1)$ and reject the hypotheses with the $j-1$ smallest $p$-values, rejecting all of them when there is no such index. Then the family-wise error rate is at most $\alpha$, whatever the dependence between the $p$-values.
\end{theorem}

\begin{proof}
If $m_0=0$, then no true null can be rejected. Otherwise consider the event $\cE$ that every true null has $p_i>\alpha/m_0$; by the union bound its probability is at least $1-m_0\cdot\alpha/m_0=1-\alpha$. On $\cE$, let $j_0$ be the first position in the ordering that holds a true null. The positions before it hold false nulls only, so $j_0\le m-m_0+1$, that is $m-j_0+1\ge m_0$, and therefore
\[
p_{(j_0)}>\alpha/m_0\ge\alpha/(m-j_0+1).
\]
The procedure stops at or before position $j_0$ and rejects no true null. The family-wise error rate is at most $\Prob(\cE^c)\le\alpha$.
\end{proof}

\begin{theorem}[Benjamini and Hochberg]\label{thm:toolkit:bh}\index{Benjamini--Hochberg procedure}
Let $R=\max\{k\ge1:p_{(k)}\le\alpha k/m\}$, or $R=0$ when no such $k$ exists, and reject the $R$ hypotheses with the smallest $p$-values. If each true-null $p$-value is valid and independent of all the other $p$-values, then the false discovery rate is at most $m_0\alpha/m\le\alpha$.
\end{theorem}

\begin{proof}
The plan is to write the false discovery rate as a sum over true nulls, and to decouple each true null's $p$-value from the number of rejections by setting it to zero.

Write $Q(k)=\#\{j:p_j\le\alpha k/m\}$. Since $p_{(k)}\le\alpha k/m$ exactly when $Q(k)\ge k$, we have $R=\max\{k:Q(k)\ge k\}$. A hypothesis $H_i$ is rejected exactly when $p_i\le\alpha R/m$: the rejected $p$-values are at most $p_{(R)}\le\alpha R/m$, and a $p$-value $p_i\le\alpha R/m$ outside the $R$ smallest would give $p_{(R+1)}\le\alpha(R+1)/m$, against the maximality of $R$. Splitting the ratio $V/\max(R,1)$ according to the true nulls rejected and the value of $R$ therefore gives
\begin{equation}\label{eq:toolkit:fdr}
\E\Bigl[\frac{V}{\max(R,1)}\Bigr]=\sum_{i\in I_0}\sum_{k=1}^m\frac1k\,\Prob\bigl(p_i\le\alpha k/m,\ R=k\bigr).
\end{equation}

Fix $i\in I_0$ and let $R^{(i)}$ and $Q^{(i)}$ be computed with $p_i$ replaced by $0$; then $Q^{(i)}(k)=Q(k)+\ind\{p_i>\alpha k/m\}$, and $R^{(i)}\ge1$ always. We claim that on the event $\{p_i\le\alpha k/m\}$ the equalities $R=k$ and $R^{(i)}=k$ are equivalent. On that event $Q^{(i)}(k')=Q(k')$ for every $k'\ge k$. If $R=k$, then $Q^{(i)}(k)=Q(k)\ge k$, and for $k'>k$, $Q^{(i)}(k')=Q(k')<k'$, so $R^{(i)}=k$. If $R^{(i)}=k$, then $Q(k)=Q^{(i)}(k)\ge k$, and for $k'>k$, $Q(k')=Q^{(i)}(k')<k'$, so $R=k$.

The variable $R^{(i)}$ depends only on the $p$-values other than $p_i$, which are independent of $p_i$. By the claim, this independence, and validity,
\[
\Prob\bigl(p_i\le\alpha k/m,\ R=k\bigr)=\Prob\bigl(p_i\le\alpha k/m\bigr)\,\Prob\bigl(R^{(i)}=k\bigr)\le\frac{\alpha k}m\,\Prob\bigl(R^{(i)}=k\bigr).
\]
Substituting into \eqref{eq:toolkit:fdr}, the inner sum is at most $\frac\alpha m\sum_k\Prob(R^{(i)}=k)=\alpha/m$, and the outer sum over the $m_0$ true nulls gives $m_0\alpha/m$.
\end{proof}

\Textcite{BenjaminiYekutieli2001} proved the same bound $m_0\alpha/m$ under a form of positive dependence, positive regression dependence of the $p$-values on those of the true nulls, and showed that under arbitrary dependence the procedure run at level $\alpha/\sum_{k=1}^m k^{-1}$ controls the false discovery rate at $\alpha$. The original papers are \textcite{Holm1979} and \textcite{BenjaminiHochberg1995}.

\section{The best of several draws}\label{sec:toolkit:orderstat}

\begin{proposition}[Order statistics]\label{thm:toolkit:orderstat}\index{order statistics}\index{expected maximum}
Let $Y_1,\dots,Y_k$ be independent with common distribution function $F$ and $\E\abs{Y_1}<\infty$, and let $M_k=\max_iY_i$.
\begin{enumerate}
\item $\Prob(M_k\le x)=F(x)^k$; when $F$ is continuous, the $q$-quantile of $M_k$ is the $q^{1/k}$-quantile of $F$.
\item $\E M_k=\int_0^\infty\bigl(1-F(x)^k\bigr)\dd x-\int_{-\infty}^0F(x)^k\dd x$.
\item For standard normal draws, $e_k=\E M_k$ satisfies $e_k\le\sqrt{2\ln k}$ for every $k\ge1$. For normal draws with mean $\mu$ and standard deviation $\sigma$, $\E M_k=\mu+\sigma e_k$.
\item For standard normal draws, put $b_k=\sqrt{2\ln k}-(\ln\ln k+\ln4\pi)/(2\sqrt{2\ln k})$. As $k\to\infty$, $\sqrt{2\ln k}\,(M_k-b_k)$ converges in distribution to the \term{Gumbel distribution}\index{Gumbel distribution}, whose distribution function is $\exp(-e^{-x})$.
\end{enumerate}
\end{proposition}

\begin{proof}
Part (i) is independence: $\{M_k\le x\}$ is the intersection of the $k$ events $\{Y_i\le x\}$. If $F$ is continuous, then $F(x)^k=q$ exactly when $F(x)=q^{1/k}$. Part (ii) is the identity $\E Z=\int_0^\infty\Prob(Z>x)\dd x-\int_{-\infty}^0\Prob(Z\le x)\dd x$, valid for every integrable $Z$, applied to $Z=M_k$, which is integrable because $\abs{M_k}\le\sum_i\abs{Y_i}$. For (iii) the case $k=1$ is $e_1=0$. For $k\ge2$ and $t>0$, Jensen's inequality and the bound of a maximum of positive numbers by their sum give
\[
e^{te_k}\le\E e^{tM_k}\le\sum_{i=1}^k\E e^{tY_i}=k\,e^{t^2/2},
\]
so $e_k\le(\ln k)/t+t/2$; the choice $t=\sqrt{2\ln k}$ gives $e_k\le\sqrt{2\ln k}$. For draws with mean $\mu$ and standard deviation $\sigma>0$ the variables $(Y_i-\mu)/\sigma$ are standard normal and $M_k=\mu+\sigma\max_i(Y_i-\mu)/\sigma$, so $\E M_k=\mu+\sigma e_k$; for $\sigma=0$ both sides equal $\mu$. Part (iv) is Theorem~1.5.3 of \textcite{LeadbetterLindgrenRootzen1983}, and we do not reproduce its proof.
\end{proof}

The bound (iii) is tight to leading order, $e_k/\sqrt{2\ln k}\to1$, but the convergence is slow, and for the small $k$ of practice the exact values of $e_k$ should be computed from (ii) by numerical integration: $e_5=1.163$, $e_{10}=1.539$, and $e_{20}=1.867$ against $\sqrt{2\ln20}=2.448$. The Gumbel law is also the limit for exponential and many other light-tailed draws, but not for all: draws with polynomial tails have the Fréchet limit, and draws bounded above, such as uniform ones, typically the Weibull limit. By the extremal types theorem these three are the only possible limits of a centered and scaled maximum \parencite[Chapter~1]{LeadbetterLindgrenRootzen1983}. \Textcite{DavidNagaraja2003} treat order statistics in general.

\section{The delta method}\label{sec:toolkit:delta}

\begin{theorem}[Delta method]\label{thm:toolkit:delta}\index{delta method}
Let $\hat\theta_n$ be estimators of $\theta\in\R^p$ with $\sqrt n(\hat\theta_n-\theta)$ converging in distribution to $\cN(0,\Sigma)$, and let $g:\R^p\to\R$ be differentiable at $\theta$. Then $\sqrt n\bigl(g(\hat\theta_n)-g(\theta)\bigr)$ converges in distribution to $\cN\bigl(0,\nabla g(\theta)^\top\Sigma\,\nabla g(\theta)\bigr)$.
\end{theorem}

\begin{proof}
The plan is to expand $g$ to first order at $\theta$ and to show that the remainder stays negligible after multiplication by $\sqrt n$. Differentiability at $\theta$ means that $g(\theta+h)=g(\theta)+\nabla g(\theta)^\top h+\norm h\,\rho(h)$ for a function $\rho$ with $\rho(0)=0$ that is continuous at $0$. With $h_n=\hat\theta_n-\theta$,
\[
\sqrt n\bigl(g(\hat\theta_n)-g(\theta)\bigr)=\nabla g(\theta)^\top\sqrt n\,h_n+\norm{\sqrt n\,h_n}\,\rho(h_n).
\]
The first term converges in distribution to $\cN(0,\nabla g^\top\Sigma\nabla g)$ by the continuous mapping theorem. In the second, $\norm{\sqrt nh_n}$ is bounded in probability because $\sqrt nh_n$ converges in distribution, and $\rho(h_n)\to0$ in probability because $h_n=(\sqrt nh_n)/\sqrt n\to0$ in probability and $\rho$ is continuous at $0$. Their product tends to $0$ in probability, and Slutsky's lemma completes the proof.
\end{proof}

For example, the compute-optimal allocation exponent of \cref{ch:fitting} is $a=\alpha_D/(\alpha_N+\alpha_D)$, and $b=1-a$. For positive exponents, $a$ decreases in $\alpha_N$ and increases in $\alpha_D$: its gradient in $(\alpha_N,\alpha_D)$ is $(-\alpha_D,\alpha_N)/(\alpha_N+\alpha_D)^2=ab\,(-1/\alpha_N,1/\alpha_D)$. Let the fitted exponents have covariance $\Sigma$, coefficients of variation $\omega_N=\Sigma_{NN}^{1/2}/\alpha_N$ and $\omega_D=\Sigma_{DD}^{1/2}/\alpha_D$, and correlation $\varrho$. The delta method gives
\[
\Var\hat a\approx a^2b^2\bigl(\omega_N^2-2\varrho\,\omega_N\omega_D+\omega_D^2\bigr)=a^2b^2\bigl[(\omega_N-\omega_D)^2+2(1-\varrho)\,\omega_N\omega_D\bigr].
\]
Because the two components of the gradient have opposite signs, a positive correlation lowers this variance below its value for uncorrelated exponents of the same precision: an error that scales both exponents by the same factor leaves $a$ unchanged. The cancellation is complete only when $\varrho=1$ and the coefficients of variation are equal, and no correlation offsets unequal ones. In particular the variance falls below both $a^2b^2\omega_N^2$ and $a^2b^2\omega_D^2$, the values it would take with one of the two exponents known without error, precisely when $\varrho>\frac12\max(\omega_N/\omega_D,\,\omega_D/\omega_N)$, which requires $\varrho>1/2$. For a quantity that increases in both exponents the gradient components would share a sign, and a positive correlation would raise its variance instead.

\section{Nonlinear least squares}\label{sec:toolkit:nls}

In a nonlinear regression $y_i=f(x_i;\beta)+\varepsilon_i$, $i=1,\dots,n$, with parameters $\beta\in\R^p$, the least-squares estimate $\hat\beta$ minimizes $\sum_i(y_i-f(x_i;\beta))^2$. Let $J$ be the $n\times p$ Jacobian with entries $\partial f(x_i;\beta)/\partial\beta_j$ at the true $\beta$.

\begin{theorem}[Asymptotics of nonlinear least squares]\label{thm:toolkit:nls}\index{nonlinear least squares}
Suppose that the errors $\varepsilon_i$ are independent and identically distributed with mean zero and variance $\sigma^2$, that $\beta$ lies in the interior of a compact parameter set, that $f(x;\cdot)$ is twice continuously differentiable, and that the regularity conditions of \textcite{Jennrich1969} hold: averages over $i$ of products of $f$ and its first and second derivatives converge uniformly in the parameters, the limit of $\frac1n\sum_i\bigl(f(x_i;\beta')-f(x_i;\beta)\bigr)^2$ is positive for every $\beta'\ne\beta$, and $J^\top J/n$ converges to a positive definite matrix $V$. Then $\hat\beta$ is consistent and $\sqrt n(\hat\beta-\beta)$ converges in distribution to $\cN(0,\sigma^2V^{-1})$. In practice the covariance of $\hat\beta$ is estimated by $\hat\sigma^2(\hat J^\top\hat J)^{-1}$, with $\hat\sigma^2$ the residual sum of squares divided by $n-p$ and $\hat J$ the Jacobian at $\hat\beta$.
\end{theorem}

\begin{proofsketch}
Expanding the normal equations $\hat J^\top(y-f(\hat\beta))=0$ to first order around $\beta$ gives $\hat\beta-\beta\approx(J^\top J)^{-1}J^\top\varepsilon$, a linear function of the errors to which the central limit theorem applies; the regularity conditions control the remainder. \Textcite{Jennrich1969} gives the proof, and \textcite{SeberWild1989} the practice.
\end{proofsketch}

Three practical consequences are used in \cref{ch:fitting}. First, if $J^\top J$ is nearly singular, then some combination of parameters is barely identified and its variance is large; the condition number of $J^\top J$, with the parameters on comparable scales, measures this. Second, the objective is in general not convex, so a fit should be started from a grid of initial values and checked for convergence. Third, with a robust loss such as Huber's \parencite{Huber1964}, which is quadratic for small residuals and linear for large ones, the covariance becomes the sandwich $A^{-1}BA^{-1}$, with $A$ the expected Hessian of the loss and $B$ the covariance of its gradient \parencite[Chapter~5]{vanderVaart1998}, and the formula above no longer applies. When in doubt, bootstrap the runs (\cref{sec:toolkit:resampling}).

\section{Variance components and designed experiments}\label{sec:toolkit:design}

In the \term{one-way random-effects model}\index{random effects}\index{variance components}, $Y_{ij}=\mu+a_i+e_{ij}$ for $I$ groups and $J$ runs per group, with all $a_i$ and $e_{ij}$ independent, the $a_i$ of mean zero and variance $\sigma_a^2$, and the $e_{ij}$ of mean zero and variance $\sigma_e^2$. A group may be an initialization seed with $J$ data orders, or a hardware configuration with $J$ seeds. The mean squares are
\[
\mathrm{MS}_A=\frac{J}{I-1}\sum_{i=1}^I(\bar Y_{i\cdot}-\bar Y_{\cdot\cdot})^2,\qquad
\mathrm{MS}_E=\frac1{I(J-1)}\sum_{i=1}^I\sum_{j=1}^J(Y_{ij}-\bar Y_{i\cdot})^2.
\]

\begin{theorem}[Variance components]\label{thm:toolkit:anova}
Consider the one-way random-effects model with $I,J\ge2$.
\begin{enumerate}
\item $\E\,\mathrm{MS}_E=\sigma_e^2$ and $\E\,\mathrm{MS}_A=\sigma_e^2+J\sigma_a^2$. Hence $\hat\sigma_e^2=\mathrm{MS}_E$ and $\hat\sigma_a^2=(\mathrm{MS}_A-\mathrm{MS}_E)/J$ are unbiased.
\item If, in addition, the $a_i$ and the $e_{ij}$ are normal and $\sigma_e^2>0$, then $I(J-1)\mathrm{MS}_E/\sigma_e^2$ and $(I-1)\mathrm{MS}_A/(\sigma_e^2+J\sigma_a^2)$ are independent and have the $\chi^2$ distributions with $I(J-1)$ and $I-1$ degrees of freedom. Consequently $(\mathrm{MS}_A/\mathrm{MS}_E)/(1+J\sigma_a^2/\sigma_e^2)$ has the $F$ distribution with $I-1$ and $I(J-1)$ degrees of freedom.
\end{enumerate}
\end{theorem}

\begin{proof}
The plan is to treat the within-group deviations and the group means separately, as two samples to which the facts about sample variances apply. For (i) the fact is this: for independent $Z_1,\dots,Z_r$ with a common mean and a common variance $v$, $\E\sum_l(Z_l-\bar Z)^2=(r-1)v$. Within group $i$, $Y_{ij}-\bar Y_{i\cdot}=e_{ij}-\bar e_{i\cdot}$, so the inner sum in $\mathrm{MS}_E$ is the sum of squared deviations of $J$ independent errors and has expectation $(J-1)\sigma_e^2$; summing over the $I$ groups and dividing by $I(J-1)$ gives $\sigma_e^2$. The group means $\bar Y_{i\cdot}=\mu+a_i+\bar e_{i\cdot}$ are independent with common mean $\mu$ and common variance $\sigma_a^2+\sigma_e^2/J$, so $\E\sum_i(\bar Y_{i\cdot}-\bar Y_{\cdot\cdot})^2=(I-1)(\sigma_a^2+\sigma_e^2/J)$, and multiplying by $J/(I-1)$ gives $J\sigma_a^2+\sigma_e^2$.

For (ii) the facts are those of normal samples used in \cref{thm:toolkit:tinterval}: the sample mean is independent of the deviations from it, and the sum of squared deviations of $r$ independent normals with a common mean and variance $v$, divided by $v$, has the $\chi^2$ distribution with $r-1$ degrees of freedom. Within group $i$ they make $\sum_j(e_{ij}-\bar e_{i\cdot})^2/\sigma_e^2$ a $\chi^2$ variable with $J-1$ degrees of freedom, independent of $\bar e_{i\cdot}$; the groups are independent, so $I(J-1)\mathrm{MS}_E/\sigma_e^2$, the sum of these $I$ variables, has the $\chi^2$ distribution with $I(J-1)$ degrees of freedom. The group means are independent normals of variance $\sigma_a^2+\sigma_e^2/J$, so $\sum_i(\bar Y_{i\cdot}-\bar Y_{\cdot\cdot})^2/(\sigma_a^2+\sigma_e^2/J)$, which equals $(I-1)\mathrm{MS}_A/(\sigma_e^2+J\sigma_a^2)$, has the $\chi^2$ distribution with $I-1$ degrees of freedom. The first variable depends only on the within-group deviations of the errors, the second only on the $a_i$ and the $\bar e_{i\cdot}$, which are independent of those deviations, so the two are independent. The ratio of two independent $\chi^2$ variables, each divided by its degrees of freedom, has the $F$ distribution by definition, and this ratio is the stated one.
\end{proof}

Under $\sigma_a^2=0$ the ratio $\mathrm{MS}_A/\mathrm{MS}_E$ itself has this $F$ distribution, which gives the $F$-test of $\sigma_a^2=0$, and solving (ii) for $\sigma_a^2/\sigma_e^2$ gives exact intervals for it; part (i) needs no normality. The \term{intraclass correlation}\index{intraclass correlation} $\sigma_a^2/(\sigma_a^2+\sigma_e^2)$ is the correlation between two runs of the same group, and it is estimated by $\hat\sigma_a^2/(\hat\sigma_a^2+\hat\sigma_e^2)$.

The estimate $\hat\sigma_a^2$ is negative whenever $\mathrm{MS}_A<\mathrm{MS}_E$, which happens often when the true component is small. It should then be reported as computed, as \cref{ch:seeds} recommends, since truncating it at zero biases it upward. Restricted maximum likelihood, the usual alternative, keeps the estimate nonnegative, but in this balanced design it does so by that very truncation: when $\mathrm{MS}_A<\mathrm{MS}_E$ it sets $\hat\sigma_a^2=0$ and estimates $\sigma_e^2$ by the total sum of squares $\sum_{i,j}(Y_{ij}-\bar Y_{\cdot\cdot})^2$ divided by $IJ-1$ \parencite[Chapter~3]{SearleCasellaMcCulloch1992}. The crossed two-way design used for seeds in \cref{ch:seeds} is analyzed in the same way.

In a \term{two-level factorial design}\index{factorial design}, each of $k$ factors is set to a low level, coded $-1$, or a high level, coded $+1$, and each of the $2^k$ combinations of levels, the \emph{cells} of the design, is run $r$ times, so that there are $r2^k$ runs. For a nonempty set $S$ of factors and levels $x=(x_1,\dots,x_k)$, put $c_S(x)=\prod_{j\in S}x_j$. The \term{effect} of $S$ is estimated by the mean outcome over the runs with $c_S=+1$ minus the mean over the runs with $c_S=-1$: a main effect when $S$ has one factor, an interaction when it has more \parencite{BoxHunterHunter2005}.

\begin{proposition}[Factorial contrasts]\label{thm:toolkit:factorial}
If the runs of a full two-level factorial design are independent with common variance $\sigma^2$, then the $2^k-1$ effect estimates are uncorrelated and each has variance $4\sigma^2/(r2^k)$, four times $\sigma^2$ divided by the number of runs.
\end{proposition}

\begin{proof}
The plan is to show that every sign $c_U$ with $U$ nonempty balances over the design, and to read off variances and covariances. Let $U$ be a nonempty set of factors and fix $j\in U$. Changing the sign of $x_j$ maps the cells one to one onto themselves and changes the sign of $c_U$, so $c_U=+1$ on exactly half of the cells, and hence on exactly half of the runs. Index the runs by $l=1,\dots,r2^k$, with levels $x^{(l)}$ and outcomes $Y_l$. The estimate of the effect of $S$ is therefore $\hat\theta_S=\frac2{r2^k}\sum_lc_S(x^{(l)})Y_l$, and its variance is $\frac4{(r2^k)^2}\sum_lc_S(x^{(l)})^2\sigma^2=4\sigma^2/(r2^k)$. For two different sets $S$ and $T$, $c_S c_T=c_{S\triangle T}$ because $x_j^2=1$, and the symmetric difference $S\triangle T$ is nonempty, so $c_{S\triangle T}$ sums to zero over the runs:
\[
\Cov(\hat\theta_S,\hat\theta_T)=\frac{4\sigma^2}{(r2^k)^2}\sum_lc_{S\triangle T}(x^{(l)})=0.\qedhere
\]
\end{proof}

Two further principles of design, set out by \textcite{Fisher1935}, complete the toolkit. \emph{Randomize} the order and placement of runs (which machine, which time, which batch of the queue), so that drifts and hidden nuisance variables cannot line up with the treatment. \emph{Block} on known nuisance variables, comparing treatments within blocks: running both arms with the same seeds, or evaluating both on the same test items, is blocking, and its analysis is the paired one of \cref{sec:toolkit:mean}. \Textcite{Montgomery2017} is a standard textbook.

\section{Sample sizes and power}\label{sec:toolkit:power}

\begin{proposition}[Runs needed for a comparison]\label{thm:toolkit:power}\index{power}\index{sample size}
Two arms have independent normal outcomes with common known standard deviation $\sigma$ and $n$ runs each, and their means differ by $\Delta$; put $\delta=\abs\Delta\sqrt{n/2}/\sigma$. The two-sided level-$\alpha$ test that rejects $\Delta=0$ when $\abs{Z}>z_{1-\alpha/2}$, with $Z=(\bar Y_1-\bar Y_2)/(\sigma\sqrt{2/n})$, has power
\[
\Phi(\delta-z_{1-\alpha/2})+\Phi(-\delta-z_{1-\alpha/2}).
\]
In particular, the power is at least $\gamma$ when $n\ge2\sigma^2(z_{1-\alpha/2}+z_\gamma)^2/\Delta^2$.
\end{proposition}

\begin{proof}
The statistic $Z$ is normal with variance $1$ and mean $\Delta\sqrt{n/2}/\sigma$, so
$\Prob(\abs Z>z_{1-\alpha/2})=\Phi(\delta-z_{1-\alpha/2})+\Phi(-\delta-z_{1-\alpha/2})$, an expression that depends on $\Delta$ only through $\delta$. The second term is nonnegative, so the power is at least $\Phi(\delta-z_{1-\alpha/2})$, which is at least $\gamma$ when $\delta\ge z_{1-\alpha/2}+z_\gamma$; solving for $n$ gives the stated bound.
\end{proof}

For example, to detect a gain of $\Delta=0.3$ points with a run-to-run standard deviation of $\sigma=0.2$ points at level $0.05$ with power $0.8$ ($z_{0.975}\approx1.960$, $z_{0.8}\approx0.842$), $n\ge2\cdot0.04\cdot2.802^2/0.09\approx6.98$, so seven runs per arm suffice. With $\sigma$ estimated from the runs the test is the two-sample $t$-test, whose power follows from the noncentral $t$ distribution with $2n-2$ degrees of freedom and noncentrality $\delta$; in the example it needs nine runs per arm, where its power is $0.85$, against $0.797$ with eight. With paired runs, $2\sigma^2$ is replaced by $\Var(Y_1-Y_2)$. \Textcite{Cohen1988} is the classic reference on power.

\section{Rank statistics}\label{sec:toolkit:ranks}

For $n\ge2$ pairs $(X_i,Y_i)$, \term{Kendall's $\tau$}\index{Kendall's tau} is
\[
\hat\tau=\binom n2^{-1}\sum_{i<j}\sgn(X_i-X_j)\,\sgn(Y_i-Y_j)\in[-1,1].
\]
A pair of indices is \term{concordant} when its term is $+1$ and discordant when it is $-1$, so that without ties $\hat\tau=(n_{\mathrm c}-n_{\mathrm d})/\binom n2$ with $n_{\mathrm c}$ concordant and $n_{\mathrm d}$ discordant pairs; ties count as zero. The coefficient depends on the data only through their ranks, so it is unchanged by strictly increasing transformations of either variable. When the pairs are independent draws of a pair $(X,Y)$, each term has expectation $\tau=\Prob\bigl((X_1-X_2)(Y_1-Y_2)>0\bigr)-\Prob\bigl((X_1-X_2)(Y_1-Y_2)<0\bigr)$, for two independent draws $(X_1,Y_1),(X_2,Y_2)$; this is the population value, and $\hat\tau$ is unbiased for it.

\begin{proposition}[Kendall's $\tau$ under independence]\label{thm:toolkit:kendall}
If the pairs $(X_i,Y_i)$ are independent draws of a pair $(X,Y)$ in which $X$ and $Y$ are independent with continuous distributions, then $\E\hat\tau=0$ and $\Var\hat\tau=2(2n+5)/(9n(n-1))$.
\end{proposition}

\begin{proof}
Write $s_{ij}=\sgn(X_i-X_j)\sgn(Y_i-Y_j)$ for $i\ne j$. Ties have probability zero, and $X_i-X_j$ and $Y_i-Y_j$ are independent and symmetric about zero, so $\E s_{ij}=0$, which gives $\E\hat\tau=0$. The variance of $\sum_{i<j}s_{ij}$ is the sum of the covariances of all ordered pairs of its terms. Each term contributes $\E s_{ij}^2=1$ with itself, and terms with no index in common are independent and contribute zero. Two terms that share one index, say $s_{ij}$ and $s_{il}$ with $j\ne l$, have $\E s_{ij}s_{il}=\E[\sgn(X_i-X_j)\sgn(X_i-X_l)]\,\E[\sgn(Y_i-Y_j)\sgn(Y_i-Y_l)]=\frac13\cdot\frac13$: two such signs agree exactly when $X_i$ (or $Y_i$) is the largest or the smallest of three exchangeable continuous draws, which has probability $2/3$, so each factor is $\frac23-\frac13$. Each of the $\binom n2$ terms shares exactly one index with $2(n-2)$ others, so $\Var\sum_{i<j}s_{ij}=\binom n2\bigl(1+2(n-2)/9\bigr)=\binom n2(2n+5)/9$, and dividing by $\binom n2^2$ gives the stated variance \parencite{Kendall1938,Kendall1970}.
\end{proof}

When a sweep crosses several hyperparameters, \textcite{JiangEtAl2020} average Kendall's $\tau$ over groups of runs that differ in one hyperparameter only, and then over the hyperparameters. A correlation within such a group cannot come from the hyperparameters that the group holds fixed; \cref{ch:sharpness} studies this \emph{granulated} coefficient.

The \term{probability of improvement}\index{probability of improvement} of arm~1 over arm~2 is $\theta=\Prob(Y_1>Y_2)+\frac12\Prob(Y_1=Y_2)$ for independent outcomes $Y_1$ and $Y_2$ of the two arms.

\begin{proposition}[Mann--Whitney estimate]\label{thm:toolkit:mannwhitney}
For independent samples $Y_{11},\dots,Y_{1n_1}$ of arm~1 and $Y_{21},\dots,Y_{2n_2}$ of arm~2,
\[
\hat\theta=\frac1{n_1n_2}\sum_{i=1}^{n_1}\sum_{j=1}^{n_2}\Bigl(\ind\{Y_{1i}>Y_{2j}\}+\tfrac12\ind\{Y_{1i}=Y_{2j}\}\Bigr)
\]
is unbiased for $\theta$. If the two arms have the same continuous distribution, then $\theta=1/2$ and $\Var\hat\theta=(n_1+n_2+1)/(12n_1n_2)$.
\end{proposition}

\begin{proof}
Each pair $(Y_{1i},Y_{2j})$ has the distribution of $(Y_1,Y_2)$, so each summand has expectation $\theta$, and so has their average. Under a common continuous distribution, $Y_1$ and $Y_2$ are exchangeable and tie with probability zero, so $\theta=1/2$ and each summand $\phi_{ij}=\ind\{Y_{1i}>Y_{2j}\}$ has variance $1/4$. Summands with no run in common are independent. Two that share the run of arm~1, $\phi_{ij}$ and $\phi_{il}$ with $j\ne l$, have covariance $\Prob(Y_{1i}>Y_{2j},\,Y_{1i}>Y_{2l})-\frac14=\frac13-\frac14=\frac1{12}$, since $Y_{1i}$ is the largest of three exchangeable continuous draws with probability $1/3$, and two that share the run of arm~2 have covariance $1/12$ for the same reason. Each summand shares its run of arm~1 with $n_2-1$ others and its run of arm~2 with $n_1-1$ others, so the variance of $\sum_{i,j}\phi_{ij}$ is $n_1n_2\bigl(\frac14+\frac{n_1+n_2-2}{12}\bigr)=n_1n_2(n_1+n_2+1)/12$; dividing by $(n_1n_2)^2$ gives the stated variance. This is the variance of the Mann--Whitney statistic $U=n_1n_2\hat\theta$ computed by \textcite{MannWhitney1947}.
\end{proof}

\section{Proportions near zero}\label{sec:toolkit:binomial}

Let $X\sim\Bin(m,p)$ count the successes on $m$ independent test items. For $x<m$ let $p_U(x)$ be the solution $q$ of $\Prob_q(X\le x)=\alpha/2$, and put $p_U(m)=1$; for $x>0$ let $p_L(x)$ be the solution $q$ of $\Prob_q(X\ge x)=\alpha/2$, and put $p_L(0)=0$.

\begin{proposition}[Binomial intervals near zero]\label{thm:toolkit:clopper}\index{Clopper--Pearson interval}\index{rule of three}
\begin{enumerate}
\item The interval $[p_L(X),p_U(X)]$ of \textcite{ClopperPearson1934} contains $p$ with probability at least $1-\alpha$, for every $p$.
\item The one-sided upper bound $p_U^{\alpha}(X)$, defined as $p_U$ with $\alpha$ in place of $\alpha/2$, satisfies $\Prob_p\bigl(p\le p_U^{\alpha}(X)\bigr)\ge1-\alpha$ for every $p$. If $X=0$, then this bound is the solution $q$ of $(1-q)^m=\alpha$, which equals $1-\alpha^{1/m}$; it is at most $\ln(1/\alpha)/m$, about $3/m$ at $\alpha=0.05$.
\end{enumerate}
\end{proposition}

\begin{proof}
For (i), fix $p$. For $x<m$ the probability $\Prob_q(X\le x)$ is continuous and strictly decreasing in $q$, so $p_U(x)<p$ holds exactly when $\Prob_p(X\le x)<\alpha/2$; for $x=m$ neither holds. Let $x^\star$ be the largest $x$ with $\Prob_p(X\le x)<\alpha/2$, or $-1$ when there is none; the cumulative probabilities increase in $x$, so the event $\{p_U(X)<p\}$ is contained in $\{X\le x^\star\}$, whose probability is less than $\alpha/2$ when $x^\star\ge0$ and zero otherwise. The same argument with $\Prob_q(X\ge x)$, which increases in $q$ for $x>0$, gives $\Prob_p(p_L(X)>p)\le\alpha/2$, and the union bound gives the coverage. For (ii), the first half of this argument with $\alpha$ in place of $\alpha/2$ gives $\Prob_p\bigl(p_U^{\alpha}(X)<p\bigr)\le\alpha$. At $x=0$ the defining equation is $(1-q)^m=\alpha$; the left side decreases from $1$ to $0$ on $[0,1]$ and equals $\alpha$ at $q=1-\alpha^{1/m}=1-e^{-\ln(1/\alpha)/m}$, which is at most $\ln(1/\alpha)/m$ because $1-e^{-t}\le t$; and $\ln20\approx3.0$.
\end{proof}

The rule of three is the subject of a short paper by \textcite{HanleyLippmanHand1983}. The Clopper--Pearson interval is conservative; the interval of \textcite{Wilson1927} is shorter and its coverage is close to nominal for most $p$, and the usual normal-approximation interval $\hat p\pm z_{1-\alpha/2}\sqrt{\hat p(1-\hat p)/m}$ fails badly near $0$ and $1$, where it can even have zero width \parencite{BrownCaiDasGupta2001}.

\notesandrefs

The tools of this appendix are classical, and each has a standard source. For the $t$ distribution and the theory of tests, \textcite{CasellaBerger2002} and \textcite{LehmannRomano2005}; for the bootstrap, \textcite{Efron1979} and \textcite{EfronTibshirani1993}; for asymptotics, including the delta method and M-estimators, \textcite{vanderVaart1998}; for designed experiments, \textcite{Fisher1935}, \textcite{BoxHunterHunter2005}, and \textcite{Montgomery2017}; for variance components, \textcite{SearleCasellaMcCulloch1992}; for extremes, \textcite{LeadbetterLindgrenRootzen1983}.

\paragraph{Translating the literature.} Statistics texts write $\bar X$, $s$, and $n$ where this book writes $\bar Y$, $S$, and $n$, and call the arms of an experiment treatments or groups. Bootstrap texts write $B$ for the number of bootstrap samples, a letter this book reserves for the batch size. Machine learning papers often report ``mean $\pm$ standard deviation''; the standard deviation $S$ describes the spread of runs, while the standard error $S/\sqrt n$ describes the uncertainty of the mean, and a 95\% interval has a half-width of about two standard errors. What \textcite{DodgeEtAl2019} call expected validation performance is the expected maximum $\E M_k$ of \cref{thm:toolkit:orderstat}.

\printchapterbib
\end{refsection}
